\documentclass[11pt,a4paper]{article}
\usepackage[margin=0.8in,top=0.85in,bottom=0.85in]{geometry}
\usepackage{amsmath,amssymb,amsfonts}
\usepackage{bm}
\usepackage{booktabs}
\usepackage{enumitem}
\usepackage{fancyhdr}
\usepackage{microtype}
\usepackage{xcolor}
\usepackage[colorlinks=true,linkcolor=black,citecolor=black,urlcolor=black]{hyperref}

% Clean, airy typography with generous breathing room
\linespread{1.10}
\setlength{\parindent}{0pt}
\setlength{\parskip}{5.5pt plus 1.5pt minus 1pt}

% Header and Footer
\pagestyle{fancy}
\fancyhf{}
\renewcommand{\headrulewidth}{0pt}
\rfoot{\footnotesize Page \thepage\ of 3}
\lfoot{\footnotesize \textit{IIT Bombay --- Department of Electrical Engineering --- EE 603: Digital Signal Processing}}

\begin{document}

% ═════════════════════════════════════════════════════════════════
% IIT BOMBAY EE 603 DEPARTMENTAL HEADER
% ═════════════════════════════════════════════════════════════════
\noindent
\begin{tabular*}{\textwidth}{@{}l @{\extracolsep{\fill}} r@{}}
\textbf{\large INDIAN INSTITUTE OF TECHNOLOGY BOMBAY} & \textbf{Autumn 2026} \\
Department of Electrical Engineering & Filter Design Laboratory \\
\textbf{EE 603: Digital Signal Processing} & Reference Note 01
\end{tabular*}

\vspace{6pt}
\hrule height 1.4pt \vspace{2pt} \hrule height 0.5pt
\vspace{20pt}

\begin{center}
{\LARGE \textbf{Digital Filter Theory \& Mathematical Formulations}}\\[8pt]
{\large \textsc{Continuous Prototypes, Bilinear Mapping, and Second-Order Sections (SOS)}}
\end{center}

\vspace{18pt}

\section*{1. Theoretical Foundations \& Overview}

Linear time-invariant (LTI) digital filtering is a foundational subject in discrete-time signal processing. Classical infinite impulse response (IIR) filter synthesis translates well-understood continuous-time prototype transfer functions into numerically stable discrete-time difference equations.

In practice, the modern synthesis pipeline comprises three systematic stages:
\begin{enumerate}[leftmargin=24pt,itemsep=3pt]
\item \textbf{Analog Prototype Approximation:} Deriving an all-pole or pole-zero continuous transfer function $H_a(s)$ satisfying magnitude and phase specifications in the normalized continuous frequency domain.
\item \textbf{Continuous-to-Discrete Mapping:} Applying the bilinear transformation with non-linear tangent frequency pre-warping to preserve critical passband and stopband boundaries.
\item \textbf{Second-Order Section (SOS) Decomposition:} Factoring high-order polynomials into cascaded biquads to eliminate coefficient quantization sensitivity in real-time execution.
\end{enumerate}

\vspace{10pt}
\section*{2. Continuous-Time Prototype Filter Families}

\subsection*{2.1 Butterworth Approximation (Maximally Flat)}
The Butterworth filter is characterized by a maximally flat magnitude response in the passband. The squared-magnitude frequency response is given by:
\begin{equation}
|H(j\Omega)|^2 = \frac{1}{1 + \left(\frac{\Omega}{\Omega_c}\right)^{2N}}
\end{equation}
where $N$ denotes the filter order and $\Omega_c$ is the $-3.01\text{ dB}$ cutoff frequency. The poles of $H(s)H(-s)$ are distributed symmetrically along a circle of radius $\Omega_c$ in the complex $s$-plane. Causality and stability require selecting the $N$ left-half plane poles:
\begin{equation}
s_k = \Omega_c \exp\left(j \frac{\pi (2k + N - 1)}{2N}\right), \quad k = 1, 2, \dots, N
\end{equation}

\subsection*{2.2 Chebyshev Approximations (Equiripple Designs)}
Chebyshev approximations minimize the peak absolute error across the specified band by utilizing Chebyshev polynomials of the first kind, $T_N(x)$, generated via the recurrence relation $T_0(x) = 1$, $T_1(x) = x$, and $T_{n+1}(x) = 2x T_n(x) - T_{n-1}(x)$.

\textbf{Type I (Passband Equiripple):}
\begin{equation}
|H(j\Omega)|^2 = \frac{1}{1 + \epsilon^2 T_N^2\left(\frac{\Omega}{\Omega_p}\right)}
\end{equation}
where $\epsilon = \sqrt{10^{R_p/10} - 1}$ parametrizes passband ripple $R_p$ (dB). The poles reside on an ellipse in the complex $s$-plane.

\textbf{Type II (Inverse Chebyshev, Stopband Equiripple):}
\begin{equation}
|H(j\Omega)|^2 = \frac{1}{1 + \left[\delta^2 T_N^2\left(\frac{\Omega_s}{\Omega}\right)\right]^{-1}}
\end{equation}
The passband is maximally flat, while the stopband exhibits equiripple behavior with transmission zeros placed directly on the imaginary axis $s = \pm j \Omega_{z,k}$.

\subsection*{2.3 Elliptic (Cauer) Rational Filters}
Elliptic filters achieve optimal transition steepness for a prescribed order $N$ by permitting equiripple behavior in both passband and stopband:
\begin{equation}
|H(j\Omega)|^2 = \frac{1}{1 + \epsilon^2 R_N^2\left(\xi, \frac{\Omega}{\Omega_p}\right)}
\end{equation}
where $R_N(\xi, x)$ is the Chebyshev rational function with selectivity parameter $\xi = \Omega_s / \Omega_p > 1$.

\subsection*{2.4 Bessel-Thomson Linear Phase Filters}
Bessel filters approximate constant group delay $\tau_g(\Omega) = -d\theta/d\Omega \approx \tau_0$ using reverse Bessel polynomials $B_N(s)$:
\begin{equation}
H(s) = \frac{B_N(0)}{B_N(s/\Omega_0)}, \quad B_N(s) = \sum_{k=0}^{N} \frac{(2N - k)!}{2^{N-k} k! (N-k)!} s^k
\end{equation}
Because the phase response is linear across the passband, Bessel filters eliminate transient ringing in step and pulse responses.

\vspace{6pt}
\begin{table}[h!]
\centering
\small
\renewcommand{\arraystretch}{1.15}
\begin{tabular}{lllll}
\toprule
\textbf{Filter Prototype} & \textbf{Passband} & \textbf{Stopband} & \textbf{Transition Width} & \textbf{Group Delay} \\
\midrule
Butterworth & Maximally Flat & Monotonic & Moderate & Moderately Linear \\
Chebyshev Type I & Equiripple & Monotonic & Sharp & Non-linear near edge \\
Chebyshev Type II & Maximally Flat & Equiripple & Sharp & Moderate Distortion \\
Elliptic (Cauer) & Equiripple & Equiripple & Optimal (Sharpest) & Highly Non-linear \\
Bessel-Thomson & Monotonic & Monotonic & Gradual & Maximally Flat \\
\bottomrule
\end{tabular}
\caption{Comparison of canonical continuous-time prototype approximations.}
\end{table}

\section*{3. Bilinear Transformation \& Frequency Pre-Warping}

The bilinear transform maps the entire left-half continuous $s$-plane into the interior of the discrete unit circle $|z| < 1$ via trapezoidal integration:
\begin{equation}
s = \frac{2}{T_s} \frac{1 - z^{-1}}{1 + z^{-1}} \iff z = \frac{1 + \frac{T_s}{2} s}{1 - \frac{T_s}{2} s}
\end{equation}
where $T_s = 1/f_s$ is the sampling interval. Evaluating along $s = j\Omega$ and $z = e^{j\omega}$ yields the non-linear frequency warping relation:
\begin{equation}
\Omega = \frac{2}{T_s} \tan\left(\frac{\omega}{2}\right) \iff \omega = 2 \arctan\left(\frac{\Omega T_s}{2}\right)
\end{equation}
To ensure synthesized discrete cutoff frequencies $\omega_0$ match the engineering targets, continuous prototype frequencies must be pre-warped prior to polynomial transformation:
\begin{equation}
\Omega_{\text{pre}} = 2 f_s \tan\left(\frac{\pi f_{\text{target}}}{f_s}\right)
\end{equation}

\section*{4. Second-Order Sections (SOS) Cascading}

High-order rational transfer functions $H(z) = B(z)/A(z)$ exhibit extreme coefficient sensitivity, where slight numerical rounding causes catastrophic pole displacement outside the unit circle. Therefore, $H(z)$ is decomposed into a cascade of $K = \lceil N/2 \rceil$ second-order sections:
\begin{equation}
H(z) = g \cdot \prod_{k=1}^{K} \frac{b_{0,k} + b_{1,k} z^{-1} + b_{2,k} z^{-2}}{1 + a_{1,k} z^{-1} + a_{2,k} z^{-2}}
\end{equation}

\subsection*{4.1 Direct Form II Transposed Difference Equations}
Each biquad section is implemented using the Direct Form II Transposed topology:
\begin{align}
y_k[n] &= b_{0,k} x_k[n] + w_{1,k}[n-1] \\
w_{1,k}[n] &= b_{1,k} x_k[n] - a_{1,k} y_k[n] + w_{2,k}[n-1] \\
w_{2,k}[n] &= b_{2,k} x_k[n] - a_{2,k} y_k[n]
\end{align}
This topology requires only two delay state registers ($w_1, w_2$) per biquad and minimizes internal accumulator overflow in real-time execution.

\vspace{6pt}
\section*{5. Canonical References}
\begin{enumerate}[label={[\arabic*]},leftmargin=22pt,itemsep=2pt]
\item A.~V.~Oppenheim and R.~W.~Schafer, \textit{Discrete-Time Signal Processing}, 3rd ed., Prentice Hall, 2009.
\item J.~G.~Proakis and D.~G.~Manolakis, \textit{Digital Signal Processing: Principles, Algorithms, and Applications}, 4th ed., Pearson, 2007.
\item S.~K.~Mitra, \textit{Digital Signal Processing: A Computer-Based Approach}, 4th ed., McGraw-Hill, 2011.
\item U.~Z\"olzer, \textit{DAFX: Digital Audio Effects}, 2nd ed., John Wiley \& Sons, 2011.
\end{enumerate}

\end{document}
